\documentclass[conference]{IEEEtran}
\usepackage{cite}
\usepackage{amsmath,amssymb,amsfonts}
\usepackage{algorithmic}
\usepackage{graphicx}
\usepackage{textcomp}
\usepackage{xcolor}
\usepackage{hyperref}
\usepackage{booktabs}
\usepackage{float}

\def\BibTeX{{\rm B\kern-.05em{\sc i\kern-.025em b}\kern-.08em
    T\kern-.1667em\lower.7ex\hbox{E}\kern-.125emX}}

\begin{document}

\title{Wildfire Prediction using Enhanced Fire Weather Index (FWI) Features and Optimized Random Forest in the Indian Subcontinent}

\author{\IEEEauthorblockN{1\textsuperscript{st} Given Name Surname}
\IEEEauthorblockA{\textit{dept. name of organization (of Aff.)} \\
\textit{name of organization (of Aff.)}\\
City, Country \\
email address or ORCID}
\and
\IEEEauthorblockN{2\textsuperscript{nd} Given Name Surname}
\IEEEauthorblockA{\textit{dept. name of organization (of Aff.)} \\
\textit{name of organization (of Aff.)}\\
City, Country \\
email address or ORCID}
}

\maketitle

\begin{abstract}
Wildfires represent one of the most significant ecological and economic threats globally, with increasing frequency in the Indian subcontinent due to climatic shifts. This paper presents a robust predictive model for wildfire occurrences using an integration of ERA5-Land meteorological data and comprehensive fire records. To address the inherent complexity of wildfire triggers, we introduce 19 engineered features derived from the Canadian Fire Weather Index (FWI) components and raw meteorological variables, including complex interaction terms and ratios. A severe class imbalance in the dataset was addressed through controlled undersampling. We evaluated multiple machine learning architectures, with an optimized Random Forest model emerging as the superior performer. Hyperparameter optimization via Randomized Search yielded a model with an accuracy of 83\% and a stable mean F1-score of 0.8360 across cross-validation folds. Statistical analysis of feature importance revealed that FWI-based interaction terms significantly enhance predictive precision compared to raw meteorological data alone. Our findings demonstrate the feasibility of a high-accuracy, FWI-enhanced wildfire prediction system tailored for the diverse climatic zones of India.
\end{abstract}

\begin{IEEEkeywords}
Wildfire Prediction, Random Forest, Fire Weather Index (FWI), Hyperparameter Tuning, Machine Learning, Feature Engineering.
\end{IEEEkeywords}

\section{Introduction}
Wildfires are among the most destructive natural disasters, causing irreversible damage to ecosystems, significant economic losses, and pose a severe threat to human health through direct injury and atmospheric pollution. Globally, the rising temperatures and shifting precipitation patterns associated with climate change have created conducive conditions for more frequent and intense fire events. The Indian subcontinent, particularly regions such as the Himalayas, central deciduous forests, and the Western Ghats, has observed a worrying upward trend in wildfire incidents.

In the Indian context, wildfire management has traditionally been more reactive than proactive. Standard methods often rely on the detection of 'Active Fire Points' (AFPs) via satellite-based sensors such as MODIS (Moderate Resolution Imaging Spectroradiometer) and VIIRS (Visible Infrared Imaging Radiometer Suite). However, detection is inherently retrospective; by the time a thermal anomaly is identified, the fire is often already established, making containment difficult. Predictive modeling, which utilizes historical weather patterns and fire incidents, provides a strategic advantage. By forecasting high-risk zones several days in advance, resource allocation for firefighting and community warnings can be significantly optimized.

The Canadian Forest Fire Weather Index (FWI) System is a globally recognized method for assessing wildfire hazard. It integrates meteorological variables---temperature, relative humidity, wind speed, and 24-hour precipitation---to produce indices representing fuel moisture and fire behavior potential. However, the direct application of standard FWI values in the subtropical and tropical climates of India often results in suboptimal performance due to differing vegetation types and extreme moisture flux. This research addresses this gap by proposing an 'Enhanced FWI' framework. We hypothesize that by generating secondary interaction features---such as the ratio of FWI to absolute temperature or the combined effect of wind and dryness---we can provide machine learning models with more nuanced discriminative power.

This paper makes the following contributions:
\begin{itemize}
    \item Integration of multi-source ERA5-Land reanalysis data with local fire incidence records for high-fidelity modeling.
    \item Development of 19 engineered features that capture non-linear interactions between weather and fire risk.
    \item A comparative analysis of state-of-the-art ensemble models (Random Forest and XGBoost) tailored for wildfire classification.
    \item Detailed hyperparameter optimization using Randomized Search to ensure model stability and generalizability.
\end{itemize}

\section{Literature Review}
The field of wildfire prediction has evolved from simple statistical correlations to complex machine learning pipelines. Early models primarily focused on the correlation between soil moisture and fire occurrence. However, as dataset complexity grew, linear models were found lacking.

\subsection{Traditional Indices}
The Canadian FWI, the National Fire Danger Rating System (NFDRS) used in the USA, and the McArthur Forest Fire Danger Index (FFDI) in Australia are the pillars of traditional fire risk assessment. These indices are physically based models that calculate moisture levels in different fuel layers (duff, shallow, and deep). While effective, they often require local tuning. In India, researchers have attempted to use FWI with varying degrees of success, often finding that static thresholds do not apply across the entire country.

\subsection{Machine Learning Advancements}
Recent literature has seen a surge in the application of Machine Learning (ML) to improve these physical models. Support Vector Machines (SVM) were initially popular due to their ability to handle high-dimensional spaces, though they often suffer from long training times on large datasets. Artificial Neural Networks (ANNs) offer great flexibility but can be 'black boxes' that are difficult to interpret for disaster management officials.

Ensemble methods, such as Random Forest (RF) and Extreme Gradient Boosting (XGBoost), have emerged as the current state-of-the-art for tabular wildfire data. RF is particularly lauded for its resilience to outliers and its ability to rank feature importance, which is critical for scientific interpretability. Studies have shown that RF can effectively capture the complex interactions between humidity and temperature that typically precede a fire event.

\subsection{Feature Engineering and Data Imbalance}
A common challenge in wildfire prediction is class imbalance. Fire events are relatively rare compared to non-fire days, leading models to favor the 'No Fire' majority class, resulting in high overall accuracy but near-zero recall for actual fires. Techniques such as SMOTE (Synthetic Minority Over-sampling Technique) and controlled undersampling have been explored to calibrate model sensitivity. Furthermore, the role of feature interaction has been increasingly recognized. Modeling fire not just as a function of temperature, but as a function of the \textit{interaction} between low humidity and high wind speed, has shown significant improvements in precision.

\section{Data Sources and Features}
\subsection{Dataset Integration}
The primary source of meteorological data is the ERA5-Land reanalysis, providing high-resolution temporal data for temperature, relative humidity, wind speed, and precipitation. Fire incident data was sourced from historical records spanning the period 2004--2019. The integrated dataset, referred to as \textit{FLAGED\_data_enhanced}, contains approximately 394,560 samples.

\subsection{Detailed Feature Engineering and Selection}
A total of 30 features were utilized in the final model. These consist of 11 raw meteorological and vegetation indices and 19 engineered interaction terms. Table I provide a categorical breakdown of the features that exhibited the highest predictive power during our recursive feature elimination (RFE) phase.

\begin{table}[h]
\caption{Comprehensive Feature Description and Importance}
\centering
\begin{tabular}{@{}lll@{}}
\toprule
\textbf{Feature Group} & \textbf{Examples} & \textbf{Significance} \\ \midrule
Meteorological & Temp, Humidity, Rainfall & Baseline triggers. \\
Vegetation & LAI\_HV, LAI\_LV, CVH & Fuel density metrics. \\
FWI Components & BUI, ISI, DC & Dryness and spread potential. \\
Interaction Terms & temp\_wind\_interaction & Synergy of heat and air flow. \\
Ratios & fwi\_rainfall\_ratio & Risk vs. suppression factor. \\
Categorical & fwi\_category & Discretized risk thresholds. \\ \bottomrule
\end{tabular}
\end{table}

The inclusion of interaction terms, particularly the product of wind speed and dryness index, allows the model to differentiate between 'hot and still' days (lower risk) and 'hot and windy' days (extreme risk), which raw variables alone often fail to distinguish.

\section{Methodology}
The methodology employed in this study follows a rigorous pipeline, from data cleansing to advanced iterative optimization.

\subsection{Data Cleansing and Scaling}
The raw dataset contained some missing values in the vegetation indices (LAI) due to satellite orbital constraints. These were imputed using a rolling mean strategy to maintain temporal consistency. Following imputation, features were scaled using the \textit{StandardScaler} algorithm, which recenters the data around a mean of zero and scales to unit variance. This is essential for gradient-based models and improves the convergence speed of ensemble methods.

\subsection{Addressing Class Imbalance}
The initial dataset distribution was highly skewed, with fire incidents representing less than 2\% of total observations. To prevent model bias towards the majority class, we implemented a random undersampling strategy. The majority 'No Fire' class was sampled to exactly match the size of the 'Fire' class, resulting in a balanced training set of 13,558 total samples.

\subsection{Model Selection and Comparison}
We compared two dominant ensemble architectures:
\begin{enumerate}
    \item \textbf{Random Forest (RF)}: An averaging ensemble that reduces variance by building multiple independent decision trees.
    \item \textbf{XGBoost}: A boosting ensemble that minimizes loss by sequentially correcting errors from previous trees.
\end{enumerate}
Preliminary tests showed that RF provided more stable results with less overfitting on the interaction features, leading us to select it for final hyperparameter tuning.

\subsection{Hyperparameter Optimization}
To find the optimal configuration, we utilized a \textit{RandomizedSearchCV} approach. This method is computationally more efficient than exhaustive Grid Search, as it samples from a wide parameter space. The search focused on maximizing the F1-score to ensure high predictive performance for both classes. The search space included:
\begin{table}[h]
\caption{Hyperparameter Search Space and Results}
\centering
\begin{tabular}{@{}lll@{}}
\toprule
\textbf{Parameter} & \textbf{Search Space} & \textbf{Optimal Value} \\ \midrule
n\_estimators & 100 -- 1000 & 900 \\
max\_depth & 5 -- 30 & 20 \\
min\_samples\_split & 2 -- 20 & 20 \\
max\_features & 'sqrt', 0.5, 0.9 & 0.9 \\
criterion & 'gini', 'entropy' & 'gini' \\ \bottomrule
\end{tabular}
\end{table}

\section{Results and Discussion}
The optimized Random Forest model was evaluated on a held-out test set (20\%). 

\subsection{Performance Metrics Analysis}
The model demonstrated strong discriminative performance. The detailed classification metrics are summarized in Table III.

\begin{table}[h]
\caption{Detailed Classification Report}
\centering
\begin{tabular}{@{}llll@{}}
\toprule
\textbf{Class} & \textbf{Precision} & \textbf{Recall} & \textbf{F1-Score} \\ \midrule
No Fire & 0.86 & 0.80 & 0.82 \\
Fire & 0.81 & 0.87 & 0.84 \\ \midrule
\textbf{Accuracy} & \multicolumn{3}{c}{\textbf{0.83}} \\ \bottomrule
\end{tabular}
\end{table}

The high recall (0.87) for the 'Fire' class is particularly noteworthy, as it signifies that the model successfully captures nearly 90\% of actual fire events, a critical requirement for early warning systems.

\subsection{Visual Analysis of Results}
\begin{figure}[h]
\centering
\includegraphics[width=0.45\textwidth]{feature_importance.png}
\caption{Feature Importance Ranking for the Optimized Random Forest Model.}
\label{fig:importance}
\end{figure}

\textit{1) Feature Importance:} The Gini importance ranking (Fig. 1) confirms that interaction terms like the \textit{Fire Danger Index} play a pivotal role.

\begin{figure}[h]
\centering
\includegraphics[width=0.4\textwidth]{confusion_matrix.png}
\caption{Confusion Matrix of the Random Forest Classifier on the test set.}
\label{fig:confusion}
\end{figure}

\textit{2) Confusion Matrix:} Figure 2 illustrates the distribution of true versus predicted labels, showing minimal false negatives (missing actual fires).

\begin{figure}[h]
\centering
\includegraphics[width=0.5\textwidth]{fwi_heatmap.png}
\caption{Geospatial Heatmap of Average FWI overlaid with recorded fire incidents.}
\label{fig:heatmap}
\end{figure}

\textit{3) Spatial Heatmap:} By interpolating FWI values and fire locations, we observe a high spatial correlation between elevated FWI zones and historical ignition points (Fig. 3).

\section{Discussion on Feature Interactions}
A significant finding of this study is the superiority of the engineered ratio and interaction features. For instance, the \textit{fwi\_temp\_ratio} effectively captures situations where a moderate FWI becomes high-risk due to extreme ambient heat. Standard models typically treat temperature and FWI as independent, missing this synergistic effect.

\section{Conclusion and Future Work}
\subsection{Summary of Contributions}
This research demonstrates that wildfire prediction in the complex Indian landscape can be significantly improved by augmenting standard meteorological data with engineered Fire Weather Index (FWI) interaction terms. Our optimized Random Forest model, achieving 83\% accuracy and 0.87 recall, provides a reliable foundation for a proactive fire management system. The methodology of class-balanced training and feature interaction modeling proves essential for tackling the rare-event nature of wildfires.

\subsection{Future Directions and Scalability}
While our current model performs exceptionally well on historical reanalysis data, several avenues for future research exist:
\begin{itemize}
    \item \textbf{Deep Learning Integration}: Future studies could explore Long Short-Term Memory (LSTM) networks or Transformers to capture the temporal dependencies and 'memory' of fuel moisture over weeks leading up to fire season.
    \item \textbf{Live Satellite Feeds}: Integrating real-time soil moisture data from SMAP or thermal feeds from Sentinel-2 would allow for near-real-time risk dynamic updates.
    \item \textbf{Topographic Influence}: Incorporating digital elevation models (DEM) to account for slope and aspect would further refine spatial spread prediction, as fires tend to move faster uphill.
    \item \textbf{Edge Deployment}: Optimizing the model for deployment on low-power IoT devices in forest stations for localized, off-grid monitoring.
\end{itemize}
Ultimately, the goal is to transit from a research-scale model to a operational national wildfire dashboard that serves local forest officers across India.

\begin{thebibliography}{00}
\bibitem{b1} M. G. Wagner et al., ``The Development and Structure of the Canadian Forest Fire Weather Index System,'' \textit{Canadian Forestry Service}, 1987.
\bibitem{b2} J. Chen et al., ``Machine learning algorithms for wildfire prediction: A review,'' \textit{Science of The Total Environment}, 2021.
\end{thebibliography}

\end{document}
